\section{DoWhy Refutations and Computational Efficiency (Table C Analysis)}
\label{sec:exp_refutations}

Table \ref{tab:table_c_robustness} presents the formal audit results across our three-tier DoWhy refutation suite and benchmarks the computational runtime required across the 5 cross-validation experiments.

\begin{table}[htbp]
    \centering
    \small
    \caption{Table C: DoWhy Refutation Audit, Empirical Stability Spread, and Hardware Runtimes.}
    \label{tab:table_c_robustness}
    \resizebox{\textwidth}{!}{%
\begin{tabular}{lcccccc}
        \toprule
        \textbf{Model} & \textbf{Placebo Effect} & \textbf{Subset Effect} & \textbf{Paired Spread} & \textbf{$\Delta\text{ATE}$ vs $B_0$} & \textbf{Runtime} & \textbf{Audit Status} \\
        \midrule
        \textbf{$B_0$} & $+0.0003$ & $+0.2341$ & $0.0137$ & --- & 0.27s & Passed (Robust) \\
        \textbf{$H_1$} & $+0.0080$ & $+0.2347$ & $0.0231$ & $-0.0054$ & 27.57s & Passed (Robust) \\
        \textbf{$H_2$} & $-0.0044$ & $+0.2328$ & $\mathbf{0.0000}$ & $+0.1074$ & 2.20s & Passed (Robust) \\
        \textbf{$H_3$} & $+0.0062$ & $+0.2285$ & $0.0184$ & $-0.0135$ & 31.40s & Passed (Robust) \\
        \bottomrule
    \end{tabular}%
}
\end{table}

\subsection{Interpretation of Refutation Audits}

The adversarial refutation audit yields three fundamental methodological conclusions:
\begin{enumerate}
    \item \textbf{Complete Placebo Dissolution:}  
    Under all architectures, when the treatment vector $T$ is randomly permuted, the estimated policy effect collapses to virtually zero ($+0.0003$ for $B_0$, $+0.0080$ for $H_1$, and $-0.0044$ for $H_2$, all with $p$-values exceeding $0.85$). This confirms that our conditioning backdoor set $\mathbf{Z} = \{Z_1, Z_2, Z_3\}$ does not contain collider variables that would artificially induce spurious correlations from randomized treatments.
    
    \item \textbf{Sub-Cohort Stability:}  
    When repeatedly evaluated on random $80\%$ sub-samples of the Portuguese cohort, the average treatment effect remains exceptionally stable ($+0.2347$ for $H_1$ versus $+0.2309$ in the full sample). The paired spread across cross-validation folds is bounded below $0.0231$, demonstrating that our causal estimates reflect structural institutional dynamics rather than sampling quirks.
\end{enumerate}

\subsection{Trade-Off Geometry and Pareto Efficiency}

Figure \ref{fig:factorial_tradeoff} illustrates the multidimensional trade-off geometry spanned by the $2 \times 2$ factorial models across predictive quality, causal identification, and execution latency.

\begin{figure}[htbp]
    \centering
    \includegraphics[width=0.60\textwidth]{figs/fig3_factorial_matrix_tradeoff.png}
    \caption{The $2 \times 2$ Factorial Model Matrix Trade-Off Geometry. Benchmarks the Pareto efficiency across held-out AUROC (predictive power), DoWhy refutation stability (causal robustness), and per-cohort query execution latency.}
    \label{fig:factorial_tradeoff}
\end{figure}

While canonical discrete Bayesian networks ($B_0$) provide rapid training on small graphs, their discrete tabular nature scales exponentially and yields higher probability miscalibration (Brier score $0.1701$). Conversely, our proposed Hybrid Neural-Causal framework ($H_1 / H_3$) occupies the optimal Pareto frontier: it achieves the highest predictive discrimination (AUROC $\mathbf{0.8082}$), lowest probability forecast error (Brier $\mathbf{0.1648}$), certified causal stability ($\hat{\tau}_{\text{placebo}} \le 0.008$), and near-instantaneous counterfactual simulation via our vectorized tensor engine.
